% Target shared-project file: sections/03-wp2-inference.tex
% Evidence: upc2026-wp3-v1, WP2 source 98d6970 and WP1 02b67b2.
\section{Inference from two early observations}
We stack the two measured ball centers into a six-dimensional vector
$\mathbf q=(x_1,y_1,z_1,x_2,y_2,z_2)$ and evaluate the same quantities
$F(\theta)$ with the forward model. The nominal measurement model uses
independent Gaussian coordinate errors with standard deviation
$\sigma=0.010\,\mathrm m$ and exact observation times. Thus
\[
 \log L(\theta)=-\frac{\|\mathbf q-F(\theta)\|^2}{2\sigma^2}
                 -6\log(\sqrt{2\pi}\sigma).
\]
The known launch center is held fixed. The coordinate-error scale is a
modeling assumption. Separate sensitivity runs use rounding cells of width
$0.010\,\mathrm m$, either alone or after Gaussian error.
The rounding-after-Gaussian cell probability is evaluated by integrating
the Gaussian density over each observed cell. Tail differences are evaluated
on a logarithmic scale.

For B0, a flat initial-velocity reference prior gives an analytic Gaussian
posterior. Its mean is
\[
 \widehat{\mathbf v}_0=
 \frac{\sum_i t_i(\mathbf q_i-\mathbf r_0+
                 \tfrac12gt_i^2\mathbf e_z)}{\sum_i t_i^2},
 \qquad
 \operatorname{Cov}(\mathbf v_0\mid\mathbf q)
 =\frac{\sigma^2}{\sum_i t_i^2}I.
\]
B1 and B2 propagate jointly sampled latent frame centers under a flat
coordinate reference prior. B2 supplies spatial crossing positions without
an arrival-time estimate.

The aerodynamic candidates use bounded velocity priors and a truncated
normal drag prior. M2 and M3 also use centered normal transverse coefficients
restricted to a disk. We sample aerodynamic parameters from their prior,
construct a conditional Gaussian velocity proposal from the two-frame
residuals, and correct each draw by its full likelihood-to-proposal ratio.
The aerodynamic prior cancels its matching proposal density.
The frozen public-run diagnostics include attempted draws and effective
sample size. The original batch runner omitted proposal rejection counts;
these are restored in a separately versioned exploratory revision. Its
additional calibration diagnostics must not be attributed to the frozen run.
Reference-curve and wind sensitivities are separately refitted to the same
early observations.

Each weighted draw is propagated to its first terminal event. Goal entry,
plane arrival, ground-first and physical non-arrival probabilities use the
entire assigned weight. Unresolved numerical weight is retained and adds
to each event-probability upper bound. Crossing intervals condition on
plane arrival, while the reported entry mode conditions on the ball center
entering the goal. The mode uses a Gaussian kernel evaluated on an interior
$0.02\,\mathrm m$ grid, with $0.05\,\mathrm m$ bandwidth and a
$0.10\,\mathrm m$ sensitivity check.

The frozen validation protocol uses highest-mass $0.10\,\mathrm m$ square
cells anchored at the origin for joint regions, including cells tied at the
mass cutoff. Conditional statistics require effective support of at least
100. B0 was retained by the frozen selection rules, with B2 as a descriptive
blind comparator. All final numerical statements refer to
\texttt{upc2026-wp3-v1}.
